%%%PAGE 106 rot=0 legibility=good summary=判断双缝干涉亮暗条纹的方法（路程差条件）；薄膜干涉（楔形膜、亮暗条件）；干涉法检查平面；电磁波与机械波对比表；麦克斯韦电磁场理论三条；LC振荡要点
%%%BLOCK id=106-01 ch=15 sec=双缝干涉·亮暗条纹判断 kind=knowledge alt=none cont=none y=0.03-0.22
\textbf{判断双缝干涉亮暗条纹的方法}

%FIG p106_a 0.115 0.045 0.335 0.15
\notefig[0.4]{figures/p106_a.png}

对于两步调一致的光源 % CHECK: 原文"两步调"，疑漏"个"

$r_2-r_1=k\lambda\ (k=0,1,2)$\quad 光屏上出现亮条纹 % 原稿"2"后仅一点，疑为省略号

$r_2-r_1=(2k+1)\dfrac{\lambda}{2}\ (k=0,1,2\ldots)$\quad 光屏上出现暗条纹,

对于两步调相反的光源，产生亮暗纹与上面的正好相反
%%%END
%%%BLOCK id=106-02 ch=15 sec=薄膜干涉 kind=knowledge alt=none cont=none y=0.23-0.44
\textbf{薄膜干涉}

%FIG p106_b 0.02 0.232 0.49 0.405
\notefig[0.55]{figures/p106_b.png}

$\Delta r=n\lambda$\quad $n=1,2,3$\quad 亮条纹

$\Delta r=(2n+1)\dfrac{\lambda}{2}$\quad $(n=0,1,2\ldots)$\quad 暗条纹

由重力作用成上薄下厚的楔形
%%%END
%%%BLOCK id=106-03 ch=15 sec=干涉法检查平面 kind=knowledge alt=none cont=none y=0.44-0.58
\textbf{干涉法检查平面.}

%FIG p106_c 0.125 0.472 0.495 0.573
\notefig[0.5]{figures/p106_c.png}

\begin{itemize}
\item 若平面平整，则观察到平行等间距的明暗相间条纹 % 原稿"平面"二字连写成一个字形(似"辅")，据文意读作"平面"
\item 若平面不平整，则干涉条纹发生弯曲
\end{itemize}
%%%END
%%%BLOCK id=106-04 ch=14 sec=电磁波与机械波对比 kind=summary alt=08 cont=none y=0.60-0.80
\begin{center}
\begin{tabular}{@{}lp{0.34\linewidth}p{0.4\linewidth}@{}}
\toprule
 & 电磁波 & 机械波 \\
\midrule
产生 & 由周期性变化的电场、磁场产生. & 由质点（波源）的振动产生 \\
波的特点. & 横波 & 纵波或横波. \\
波速 & 真空中等于光速$c=3\times10^{8}\unit{m/s}$ & 在空气中不大（声波在空气中$340\unit{m/s}$） \\
是否需要介质 & 不需要介质，可在真空中传播 & 必须有介质，真空中不能传播. \\
能量传播 & 电磁能 & 机械能. \\
\bottomrule
\end{tabular}
\end{center}
%%%END
%%%BLOCK id=106-05 ch=14 sec=麦克斯韦电磁场理论 kind=knowledge alt=none cont=none y=0.80-0.89
% 原稿此三行用左大括号括在一起
\begin{itemize}
\item 均匀变化的电场（磁场）产生 均匀变化的磁场（电场） % CHECK: 原文如此(物理上应为稳定的磁场(电场))
\item 非均匀变化的电场（磁场）产生 非均匀变化的磁场（电场）.
\item 周期性振荡的电场（磁场）产生 周期性振荡的磁场（电场）.
\end{itemize}
%%%END
%%%BLOCK id=106-06 ch=14 sec=LC振荡 kind=note alt=none cont=none y=0.89-1.00
$I\propto$磁场能\qquad $U\propto$电场能.\qquad $T$=2次充放电.

\struck{$U_c$的导数为$U_L$}

\struck{互为导数} % 原稿上述两行均被横线划去
%%%END
%%%PAGEEND 106
